\section{Теорема об изоморфизме для линейных пространств. Координатное пространство}

\begin{theorem}[Об изоморфизме конечномерных линейных пространств]
Два конечномерных линейных пространства $V$ и $W$ над одним и тем же полем $F$ изоморфны тогда и только тогда, когда их размерности равны:
\[
V \cong W \quad \Longleftrightarrow \quad \dim V = \dim W.
\]
\end{theorem}

\begin{proof}
\begin{itemize}
    \item \textbf{Необходимость ($\Rightarrow$):} Пусть $V \cong W$ и $\varphi: V \to W$ --- изоморфизм. Пусть $\dim V = n$, и $E = \{e_1, \dots, e_n\}$ --- базис в $V$. По свойству изоморфизма, система $\varphi(E) = \{\varphi(e_1), \dots, \varphi(e_n)\}$ линейно независима в $W$. Покажем, что она образует базис в $W$. Возьмем произвольный $y \in W$. Так как $\varphi$ сюръективно, существует $x \in V$ такой, что $\varphi(x) = y$. Разложим $x$ по базису $E$: $x = \xi_1 e_1 + \dots + \xi_n e_n$. Тогда:
    \[
    y = \varphi(x) = \varphi(\xi_1 e_1 + \dots + \xi_n e_n) = \xi_1 \varphi(e_1) + \dots + \xi_n \varphi(e_n).
    \]
    Следовательно, $\varphi(E)$ порождает $W$ и является его базисом. Значит, $\dim W = n = \dim V$.
    
    \item \textbf{Достаточность ($\Leftarrow$):} Пусть $\dim V = \dim W = n$. Выберем базис $E = \{e_1, \dots, e_n\}$ в $V$ и базис $E' = \{e'_1, \dots, e'_n\}$ в $W$. Определим отображение $\varphi: V \to W$ следующим образом: для любого $x = \xi_1 e_1 + \dots + \xi_n e_n$ положим
    \[
    \varphi(x) = \xi_1 e'_1 + \dots + \xi_n e'_n.
    \]
    Это отображение линейно (следует из теоремы о линейности координат). Оно инъективно, так как $\varphi(x) = 0_W \Rightarrow \xi_1 = \dots = \xi_n = 0 \Rightarrow x = 0_V$. Оно сюръективно, так как любой $y \in W$ имеет вид $y = \eta_1 e'_1 + \dots + \eta_n e'_n$, и для $x = \eta_1 e_1 + \dots + \eta_n e_n$ выполняется $\varphi(x) = y$. Следовательно, $\varphi$ --- изоморфизм.
\end{itemize}
\hfill$\blacksquare$
\end{proof}

\begin{corollary}[О координатном пространстве]
Любое $n$-мерное линейное пространство $V$ над полем $F$ изоморфно координатному пространству $F^n$.
\end{corollary}

\begin{proof}
Пространство $F^n$ состоит из столбцов высоты $n$ с элементами из $F$ и имеет размерность $n$ (стандартный базис --- единичные столбцы). По доказанной теореме, так как $\dim V = \dim F^n = n$, пространства изоморфны. Изоморфизм задается сопоставлением каждому вектору $x \in V$ столбца его координат в выбранном базисе:
\[
x \mapsto X = \begin{pmatrix} \xi_1 \\ \xi_2 \\ \vdots \\ \xi_n \end{pmatrix} \in F^n.
\]
\hfill$\blacksquare$
\end{proof}
